% !PRINTED 586
\weberpage{586}
If one develops the left side in powers of $w$, the coefficients are
rationally expressible through $\lambda$ and $\sqrt{\Delta}$, and hence
belong to the field
\[
\mathfrak{L} = \mathfrak{R}(\lambda, \sqrt{\Delta}).
\]
If one develops the right side in the same way and orders it by the
powers of $w$, then $A_\nu$ appears first in the coefficient of
$w^{\nu-1}$, and indeed carrying the factor 1. Accordingly $A_\nu$ can
be determined once the earlier coefficients have already been
determined, and it follows that all these coefficients belong to the
field $\mathfrak{L}$.

Since afterwards the numerator and denominator of the fraction $A/D$
can again be freed from common factors by rational calculation, the
following theorem results:

2. \gesperrt{If in the expression (8) the numerator and denominator
$A(\boldsymbol{s})$, $D(\boldsymbol{s})$ are freed from common factors,
then these functions belong to the field $\mathfrak{L}$.}
% [source sic: the theorem cites expression (8), although A and D were introduced in (10)]

In order to represent the function $A(x)$, we take the coefficient of
the highest power of $x$ equal to 1 and investigate for which values of
$u$ the function $\boldsymbol{s}(\mu u)$ vanishes. We denote a value
$u$ for which $\boldsymbol{s}(\mu u)$ vanishes by $2\varrho/\mu$, and
then obtain as the condition for the vanishing, by (5), one of the two
congruences
\[
2\varrho \equiv 0, 1 \quad (\mathrm{mod}\ 2, \omega).
\]
Thus $\varrho$ is a number of the field $\Omega$ which may indeed be
fractional, but can have no other denominators than a divisor of $2a$.
On the other hand one may alter $\varrho$ by a multiple of $\mu$
without the value $\boldsymbol{s}(2\varrho/\mu)$ being changed, and
since only the latter is of consequence, and since the congruence
\[
\varrho \equiv \frac{\xi}{2a} \quad (\mathrm{mod}\ \mu),
\]
where $\xi$ is an integer in $\Omega$, can be satisfied by an integer
$\varrho$, we may assume $\varrho$ to be an \gesperrt{integer} if we \gesperrt{presuppose
that the first coefficient $a$ in (1) is odd and relatively prime to
$\mu$.} We hold fast to this assumption from now on.

If, then, we let $\varrho$ run through a complete system of residues
with respect to the modulus $\mu$, with exclusion of zero, we obtain
all the values of $\boldsymbol{s}$ for which $A(\boldsymbol{s})$
vanishes, in the form:
\weq{13}{\boldsymbol{s}\left(\frac{2\varrho}{\mu}\right).}
